This paper argues that the central reason deep learning often underperforms in
cross-sectional cryptocurrency prediction is not necessarily architectural
inadequacy, but objective misspecification.
When the task is fundamentally about relative ordering, training on raw-return
targets makes the loss overly sensitive to week-to-week variation in return
scale.
Replacing those targets with rank percentiles, and pairing that change with
gradient accumulation, makes the ranking problem substantially more stable and
translates that stability into out-of-sample performance.

\paragraph{Architecture and signal type.}
The main empirical conclusion is not that one architecture dominates
universally, but that different architectures are well suited to different
signal structures.
CS-Gated performs best when predictive content comes from contemporaneous
on-chain fundamentals, implying that strong cross-sectional structure can be
exploited without explicit temporal encoding.
LSTM and TFT improve materially once Trump social media and granular ETF flow
variables are added, suggesting that temporal architectures are most valuable
when the predictors themselves possess persistence, event clustering, or
sequence-level information.
At the same time, Random Forest remains a powerful benchmark in richer
high-dimensional settings, underscoring that traditional nonlinear methods are
still highly competitive.

\paragraph{Methodological implication.}
The ablation evidence shows that the decisive improvement comes less from a
novel network design than from correcting the supervision signal and making it
trainable.
This has a broader implication for financial machine learning:
reported failures of deep learning may sometimes reflect the fact that models
were optimized against targets misaligned with the actual investment
criterion, rather than the fact that neural architectures are intrinsically
ill-suited to the domain.

\paragraph{Limitations.}
Three limitations remain central.
First, the test period covers only 49 weeks, which constrains statistical
precision.
Second, the 86-asset universe is fixed using March 2026 rankings and therefore
contains survivorship bias.
Third, the strongest deep-learning results rely on multi-seed ensembling, so
their advantage over simpler methods is not the same as single-run
repeatability or computational efficiency.

\paragraph{Future work.}
Several extensions follow naturally from the present results.
The rank-normalized objective could be tested on equity or other
cross-sectional asset universes to assess external validity.
Higher-frequency order-book, microstructure, or narrative-text signals could
be used to identify more precisely which predictors genuinely require temporal
encoding.
Finally, portfolio construction could be extended beyond decile long-short
sorting toward constrained optimization or dynamic position sizing, bringing
the predictive and execution layers of the framework into tighter alignment.
